\documentclass[11pt]{amsart}
\usepackage[margin=1.15in]{geometry}
\usepackage{amssymb,hyperref,tikz,booktabs}
\newtheorem{theorem}{Theorem}
\newtheorem{lemma}[theorem]{Lemma}
\newtheorem*{ps}{Theorem PS}
\theoremstyle{definition}\newtheorem{problem}{Problem}
\theoremstyle{remark}\newtheorem*{remark}{Remark}
\title{Non-nested matchings in ordered graphs: the cases $k=3$ and $k=4$}
\author{Vamshi Jandhyala}
\date{26 November 2026 (version 2; version 1: 25 November 2026)}
\begin{document}
\begin{abstract}
Two disjoint edges of a graph on $[n]$ are nested when one lies strictly inside the other. Bar\'at, Freschi and
T\'oth conjectured that an ordered graph on $n\ge 2k$ vertices with no $k$ pairwise disjoint, pairwise non-nested
edges has at most $(k-1)n$ edges, and that the same bound holds when only the more restricted \emph{strongly}
non-nested matchings are forbidden. De Castro disproved the conjecture for every $k\ge5$. We prove both forms for
$k=3$, we prove the plain form for $k=4$, and we show that the strong form fails for $k=4$: there is an ordered graph
on $9$ vertices with $28$ edges and no strongly non-nested $4$-matching, and the excess persists for every $n\ge9$.
So the plain form of the conjecture holds exactly for $k\le4$ and the strong form exactly for $k\le3$. Both upper
bounds study the equality case of the Bar\'at--Freschi--T\'oth count and close it with the star polygons of
multitriangulations; for $k=4$ two finite configuration checks are done by computer.
\end{abstract}
\maketitle

\section{The question}
Put the points $1,2,\dots,n$ on a line and draw arcs between some pairs of them, all above the line. How many arcs can
you draw before you are forced to find three arcs with no common endpoints, none of which sits inside another? Try it
for $n=7$ before reading on. One good answer joins every pair among $1,\dots,5$ and then joins $6$ and $7$ to $1$ and
$2$, which gives $10+4=14=2n$ arcs (Figure~\ref{fig:patterns}, bottom). Any three disjoint arcs must use $6$ or $7$.
In three pairwise non-nested disjoint arcs, the order of the left ends is the order of the right ends; the arc with the
last right end ends at $6$ or $7$, so its left end is $1$ or $2$, and the three distinct left ends would all lie in
$\{1,2\}$.

\begin{figure}[ht]
\centering
\begin{tikzpicture}[x=0.5cm,y=0.5cm,v/.style={circle,fill,inner sep=1.1pt}]
\foreach \x in {0,1,2,3} \node[v] at (\x,0) {};
\draw (0,0) to[out=70,in=110] (3,0); \draw (1,0) to[out=70,in=110] (2,0);
\node at (1.5,-0.8) {\small nested};
\begin{scope}[xshift=2.6cm]
\foreach \x in {0,1,2,3} \node[v] at (\x,0) {};
\draw (0,0) to[out=70,in=110] (2,0); \draw (1,0) to[out=70,in=110] (3,0);
\node at (1.5,-0.8) {\small crossing};
\end{scope}
\begin{scope}[xshift=5.2cm]
\foreach \x in {0,1,2,3} \node[v] at (\x,0) {};
\draw (0,0) to[out=70,in=110] (1,0); \draw (2,0) to[out=70,in=110] (3,0);
\node at (1.5,-0.8) {\small separated};
\end{scope}
\end{tikzpicture}\\[4pt]
\begin{tikzpicture}[x=0.9cm,y=0.9cm,v/.style={circle,fill,inner sep=1.3pt}]
\begin{scope}
\foreach \x in {1,...,7} \node[v,label=below:{\scriptsize\x}] at (\x-1,0) {};
\foreach \a in {1,...,5} \foreach \b in {1,...,5} {\ifnum\a<\b \draw[gray] (\a-1,0) to[out=60,in=120] (\b-1,0);\fi}
\foreach \a in {1,2} \foreach \b in {6,7} \draw[thick,dashed] (\a-1,0) to[out=60,in=120] (\b-1,0);
\end{scope}
\end{tikzpicture}
\caption{The three ways two disjoint edges can sit, and the graph on $7$ vertices with $2n=14$ edges and no three
disjoint pairwise non-nested edges: $K_5$ on $\{1,\dots,5\}$ (grey) and the edges from $6,7$ to $1,2$ (dashed).}
\label{fig:patterns}
\end{figure}

Formally, an \emph{ordered graph} is a graph on $[n]=\{1,\dots,n\}$; we write an edge as $ij$ with $i<j$ and call
$j-i$ its length. Two disjoint edges $ij$ and $st$ with $i<s$ are \emph{nested} if $t<j$, \emph{crossing} if
$s<j<t$, and \emph{separated} if $j<s$. A $k$-matching is \emph{non-nested} if no two of its edges are nested. It is
\emph{strongly non-nested} if, in addition, ``equal or crossing'' is an equivalence relation on it (equivalently,
crossing is transitive on distinct edges): the matching splits into blocks
of pairwise crossing edges, and edges in different blocks are separated. For three edges, strongly non-nested means
pairwise crossing, pairwise separated, or a crossing pair together with an edge separated from both. Let $f_k(n)$ and
$f^*_k(n)$ be the largest number of edges of an ordered graph on $[n]$ with no non-nested $k$-matching, respectively
no strongly non-nested $k$-matching. Every strongly non-nested matching is non-nested, so $f_k(n)\le f_k^*(n)$.

Bar\'at, Freschi and T\'oth \cite{BFT} proved that for $n\ge2k$
\[
(k-1)n\le f_k(n)\le f^*_k(n)\le (k-1)n+\tbinom{k-1}{2}, \tag{1}
\]
and conjectured \cite[Conjecture~2.6]{BFT} that $f_k(n)=f^*_k(n)=(k-1)n$. The lower bound comes from the graph in
Figure~\ref{fig:patterns}, generalised: $K_{2k-1}$ on the first $2k-1$ vertices, and every later vertex joined to the
first $k-1$. The question sits behind a Ramsey problem of Bar\'at \cite{CC33,BGT}: an exact Tur\'an bound would improve
the known upper bound for the smallest $m$ such that every $2$-colouring of the ordered $K_m$ has a monochromatic
non-nested $k$-matching \cite[Lemma~3.3]{BFT}. De Castro \cite{dC} showed that $f_k(n)>(k-1)n$ for every $k\ge5$ and
$n\ge2k+1$, so the conjecture is false from $k=5$ on, while the cases $k=3,4$ were left open. For $k=2$ the bounds in
(1) already agree.

\begin{theorem}\label{thm:k3}
For every $n\ge6$, $f_3(n)=f^*_3(n)=2n$.
\end{theorem}

\begin{theorem}\label{thm:k4}
For every $n\ge9$, $f^*_4(n)\ge 3n+1$. In particular the strong form of \cite[Conjecture~2.6]{BFT} fails for $k=4$.
\end{theorem}

\begin{theorem}\label{thm:k4plain}
For every $n\ge8$, $f_4(n)=3n$.
\end{theorem}

Together with \cite{dC}, the plain form of the conjecture holds exactly for $k\le4$ and the strong form exactly for
$k\le3$; Theorem~\ref{thm:k4} is what separates the two forms at $k=4$. Theorem~\ref{thm:k4plain} is proved in
Section~\ref{sec:k4plain}. Its proof follows the same plan as that of Theorem~\ref{thm:k3}, with $3$-triangulations
in place of $2$-triangulations, but the equality case is far richer and two finite parts of it are checked by
computer.

\section{Proof of Theorem~\ref{thm:k3}}
By (1) it suffices to show that an ordered graph $G$ on $[n]$, $n\ge6$, with no strongly non-nested $3$-matching has
at most $2n$ edges. Call an edge \emph{short} if its length is $1$ or $2$, and \emph{long} otherwise. The proof of the
upper bound in (1) splits the edges in exactly this way, and we first recall it for $k=3$.

The route is as follows. The count (2) below gives $e(G)\le2n+1$, with equality only if $G$ has eight short edges and
its long edges are as many as Capoyleas and Pach allow. Lemma~\ref{lem:short} shows that eight short edges form two
small stars, and Lemma~\ref{lem:fence} shows that these stars confine every other edge. In the equality case the long
edges complete a $2$-triangulation, whose star polygons (Theorem PS) then force three pairwise crossing edges.

A short edge has at most one vertex strictly between its ends, so it contains no disjoint edge, and it is never one of
three pairwise crossing edges (two of them would need that same inner vertex). Let $H$ be $G$ together
with all $2n-3$ short edges. Three pairwise crossing edges form a strongly non-nested $3$-matching, so $H$ has none.
Capoyleas and Pach \cite{CP} proved that a graph on $n$ points in convex position with no three pairwise crossing
edges has at most $4n-10$ edges; drawing $[n]$ on a circle in its natural order turns crossing in our sense into
crossing chords. Writing $s$ for the number of short edges of $G$,
\[
e(G)=e(H)-(2n-3-s)\le 2n-7+s. \tag{2}
\]
In words: the long edges number at most $2n-7$, and every short edge of $G$ adds at most one to the count.

\begin{lemma}\label{lem:short}
$G$ has at most $8$ short edges. If it has exactly $8$, they are the eight edges of length at most $2$ at two vertices
$a<b$ with $b-a\ge3$, namely $a$ joined to $a-2,a-1,a+1,a+2$ and $b$ joined to $b-2,b-1,b+1,b+2$.
\end{lemma}

\begin{proof}
Three disjoint short edges are never nested, and each crosses at most one of the others (a crossing edge must use its
inner vertex), so ``equal or crossing'' is an equivalence relation on them: they form a strongly non-nested
matching. Hence the
short edges of $G$ contain no $3$-matching. The edges of a fixed length form disjoint paths (for length $2$, one on
the odd and one on the even vertices), and a union of paths with $m$ edges contains a matching of size $\lceil
m/2\rceil$. So there are at most $4$ short edges of each length, and at most $8$ in all. (Directly: among five such edges with
left ends $p_1<\dots<p_5$, those at $p_1,p_3,p_5$ are disjoint for length $1$, and one of the triples
$(p_1,p_3,p_5)$, $(p_1,p_2,p_5)$, $(p_1,p_4,p_5)$ is for length $2$.)

Suppose there are $8$, so $4$ of each length. Four length-$1$ edges with no $3$-matching form either one path
$t,t+1,\dots,t+4$ or two paths of two edges. In the first case every length-$2$ edge must meet the path, and an edge
meeting it in one vertex must use $t+1$ or $t+3$ (removing any other vertex of the path leaves two disjoint edges). This
leaves the five candidates $(t,t+2)$, $(t+1,t+3)$, $(t+2,t+4)$, $(t-1,t+1)$, $(t+3,t+5)$. Now $(t-1,t+1)$ and $(t,t+2)$
together with $(t+3,t+4)$ form a $3$-matching, and so do $(t+3,t+5)$, $(t+2,t+4)$ and $(t,t+1)$. So at most three of
the candidates occur, a contradiction.

Hence the length-$1$ edges are $(a-1,a),(a,a+1),(b-1,b),(b,b+1)$ with $a+1<b-1$. A length-$2$ edge that avoids $a$ and
$b$ can be completed to a $3$-matching by one length-$1$ edge at $a$ and one at $b$, unless it contains both
$a-1$ and $a+1$, or both $b-1$ and $b+1$. So the length-$2$ edges are among the four at $a$ and $b$ and the two chords
$(a-1,a+1)$, $(b-1,b+1)$. The chord $(a-1,a+1)$ is disjoint from $(a-2,a)$ and from $(a,a+2)$, and either pair together
with $(b,b+1)$ is a $3$-matching; so if this chord occurs, neither length-$2$ edge at $a$ does, and symmetrically at
$b$. Either chord therefore costs at least one edge, and four length-$2$ edges are possible only when both chords are
absent and all four edges at $a$ and $b$ are present.
\end{proof}

Call the eight edges of Lemma~\ref{lem:short} the \emph{short stars} at $a$ and $b$. They fence in every other edge.

\begin{lemma}\label{lem:fence}
Suppose $G$ contains the short stars at $a<b$, and let $ij\in E(G)$.
\begin{enumerate}
\item[(L)] If $j<b$ and $a\notin\{i,j\}$, then $i\le a-2$ and $j\ge a+2$.
\item[(R)] If $i>a$ and $b\notin\{i,j\}$, then $i\le b-2$ and $j\ge b+2$.
\end{enumerate}
\end{lemma}

\begin{proof}
For (L), the edge $(b,b+1)$ is disjoint from $ij$ and separated both from $ij$ and from every edge of the star at
$a$. So every star edge at $a$ that is disjoint from $ij$ must be nested with $ij$; otherwise it forms, with $ij$ and
$(b,b+1)$, a strongly non-nested $3$-matching. If $ij$ lay entirely on one side of $a$, a star edge on the other side
would be disjoint from it and separated from it. Hence $i<a<j$. If $i=a-1$ then $(a-2,a)$ crosses $ij$, and if $j=a+1$
then $(a,a+2)$ crosses $ij$; both are excluded. Part (R) is (L) for the reversed order.
\end{proof}

The last ingredient is a structural fact about maximal graphs without three pairwise crossing chords, which are
called \emph{$2$-triangulations}. A \emph{$2$-star} is the pentagram formed by five vertices $s_0,\dots,s_4$ in
cyclic order with the five chords $s_js_{j+2}$.

\begin{ps}[Pilaud and Santos {\cite[Corollary~4.2(ii)]{PS}}]
In a $2$-triangulation of the $n$-gon, $n\ge5$, every chord joining two vertices at cyclic distance $2$ belongs to a
$2$-star \textup{(}indeed to exactly one\textup{)}.
\end{ps}

We use it in the following form. Let $p,q,r$ be cyclically consecutive and let $S$ be a $2$-star containing the chord
$pr$. A chord $s_js_{j+2}$ of a $2$-star has exactly one vertex of the star strictly on one side and two on the other.
Only $q$ lies strictly between $p$ and $r$ on the short side, so the side with one star vertex is that one and the
vertex is $q$. Hence $S$ has vertices $p,q,r,u,v$ in cyclic order, and the five chords $pr,\,qu,\,rv,\,up,\,vq$
are present.

\begin{proof}[Proof of Theorem~\ref{thm:k3}]
By (2) and Lemma~\ref{lem:short}, $e(G)\le 2n+1$, and equality forces $s=8$ and $e(H)=4n-10$. Suppose equality
holds. Then $G$ contains the short stars at some $a<b$, with $3\le a$, $b\le n-2$ and $b-a\ge3$. Moreover $H$ is a
$2$-triangulation: adding any chord would break the Capoyleas--Pach bound. Every long edge of $H$ is an edge of $G$.

Apply Theorem PS to $a-1,a,a+1$: there are $u,v$ with $a-1,a,a+1,u,v$ in cyclic order and $(a-1,u)$, $(a+1,v)$ in
$H$. If $v<a-1$, then $(v,a+1)$ is long, lies in $G$, and violates (L). So $v>a+1$, and the cyclic order gives
$a+2\le u<v$. The long edge $(a-1,u)$ violates (L) unless $u\ge b$. Hence $b\le u<v$, and $(a+1,v)$ is a long edge of
$G$ with $v>b$.

Apply Theorem PS to $b-1,b,b+1$: there are $x,y$ with $b-1,b,b+1,x,y$ in cyclic order and $(x,b-1)$, $(y,b)$,
$(y,b+1)$ in $H$. If $x>b+1$, the long edge $(b-1,x)$ violates (R). So $x<y\le b-2$, and the long edge $(y,b+1)$
violates (R) unless $y\le a$. Hence $x<y\le a$.

Now $(x,b-1)$, $(y,b)$ and $(a+1,v)$ are long edges of $G$, and
\[
x<y<a+1<b-1<b<v,
\]
using $b-a\ge3$. These three edges are pairwise crossing (Figure~\ref{fig:end}), a strongly non-nested
$3$-matching. This contradiction shows $e(G)\le2n$.
\end{proof}

\begin{figure}[ht]
\centering
\begin{tikzpicture}[x=0.62cm,y=0.62cm,v/.style={circle,fill,inner sep=1.1pt}]
\foreach \x/\l in {0/x,1.2/y,3/{a\!-\!1},4/a,5/{a\!+\!1},7/{b\!-\!1},8/b,9/{b\!+\!1},11/v}
  \node[v,label=below:{\scriptsize$\l$}] at (\x,0) {};
\foreach \x in {2,6,10} \node[v] at (\x,0) {};
\foreach \c in {4,8} \foreach \d in {-2,-1,1,2} \draw[gray] (\c,0) to[out=70,in=110] (\c+\d,0);
\draw[thick] (0,0) to[out=55,in=125] (7,0);
\draw[thick] (1.2,0) to[out=55,in=125] (8,0);
\draw[thick] (5,0) to[out=55,in=125] (11,0);
\end{tikzpicture}
\caption{The end of the proof of Theorem~\ref{thm:k3}: the short stars at $a$ and $b$ (grey) and the three long edges
forced by the $2$-stars at the chords $(a-1,a+1)$ and $(b-1,b+1)$. The three thick edges pairwise cross.}
\label{fig:end}
\end{figure}

\begin{remark}
The same argument never uses a non-nested matching that is not strongly non-nested, which is why it gives both
equalities at once. For $k=4$ this is no longer possible, by Theorem~\ref{thm:k4}.
\end{remark}

\section{Proof of Theorem~\ref{thm:k4}}
Let $G_9$ be the ordered graph on $[9]$ in which $A=\{1,2,3\}$ is independent, $B=\{4,\dots,9\}$ is complete, and
\[
N(1)=N(2)=\{4,7,8,9\},\qquad N(3)=\{4,6,7,8,9\}.
\]
It has $15+4+4+5=28=3\cdot9+1$ edges. We show that every non-nested $4$-matching $M$ of $G_9$ contains edges
$e,f,g$ such that $e$ crosses $f$, $f$ crosses $g$, and $e,g$ are separated; so crossing is not transitive on $M$, and
$G_9$ has no strongly non-nested $4$-matching.

In a non-nested matching the left ends and the right ends of the edges appear in the same order. Every vertex of $A$
precedes every vertex of $B$ and $A$ is independent, so each edge of $M$ joins $A$ to $B$ or lies inside $B$. If $t$
edges lie inside $B$, then $M$ uses $4-t\le3$ vertices of $A$ and $t+4\le6$ of $B$, so $t\in\{1,2\}$.

If $t=1$, the edges are $(1,x)$, $(2,y)$, $(3,z)$, $(c,d)$ with $x<y<z<d$. The neighbourhoods force $x=4$, $y=7$, $z=8$,
$d=9$ and $c\in\{5,6\}$. Then $(1,4)$ crosses $(2,7)$, $(2,7)$ crosses $(c,9)$, and $(1,4)$ and $(c,9)$ are separated.

If $t=2$, the edges are $(a_1,b_1)$, $(a_2,b_2)$ with $a_1<a_2$ in $A$ and $(c_1,d_1)$, $(c_2,d_2)$ inside $B$ with
$c_1<c_2$. Then $b_1<b_2<d_1<d_2$ and these six vertices are all of $B$, so $b_2\le7$; also $b_2\ge6$, since $5$ has
no neighbour in $A$. If $b_2=7$, then $d_1=8$, $d_2=9$, $b_1=4$ (as $a_1\le2$) and $c_1=5$; the edges $(a_1,4)$,
$(a_2,7)$, $(5,8)$ give the pattern. If $b_2=6$, then $a_2=3$, $b_1=4$ and $c_1=5$; the edges $(a_1,4)$, $(3,6)$,
$(5,d_1)$ give the pattern.

Finally, add to an ordered graph on $[n]$ a vertex $n+1$ joined to $1,2,3$. A non-nested matching using an edge
$(i,n+1)$ has all its left ends at most $i\le3$, because $n+1$ is the last right end, so it has at most three edges.
Repeating this step from $G_9$ gives, for every $n\ge9$, a graph with $3n+1$ edges and no strongly non-nested
$4$-matching. \qed

\begin{remark}
A SAT computation gives $f^*_4(n)=3n+1,\,3n+1,\,3n+2,\,3n+2,\,3n+2$ for $n=9,\dots,13$, against the upper bound $3n+3$
from (1).
\end{remark}

\section{Proof of Theorem~\ref{thm:k4plain}}\label{sec:k4plain}
The lower bound is (1). For the upper bound let $G$ be an ordered graph on $[n]$, $n\ge8$, with no non-nested
$4$-matching. Now an edge is \emph{short} if its length is at most $3$ and \emph{long} otherwise; $S$ is the set of
short edges of $G$, $s=|S|$, and $L$ is the number of long edges. The route is the one of Section~2. A count bounds
$e(G)$ by $3n$ unless $s\ge16$ (Section~4.1). The short edges then have a rigid structure (Section~4.2). Every
remaining long edge of a $3$-triangulation completing $G$ must be present, and a single star polygon of that
triangulation, placed next to a dense short star, forces a non-nested $4$-matching (Sections~4.3--4.5). When a few long
edges may be missing, several such stars are needed and a counting argument shows they cannot all be spoiled
(Section~4.6).

\subsection{The count}
The \emph{cyclic length} of $ij$ is $\min(j-i,\,n-j+i)$, its length when $[n]$ is drawn on a circle. An edge of
cyclic length at most $3$ has at most two vertices strictly on one side, so it lies in no set of four pairwise
crossing edges. Let $H$ be $G$ together with all $3n$ edges of cyclic length at most $3$. Four pairwise crossing edges
of $G$ form a non-nested $4$-matching, so $H$ has no four pairwise crossing edges, and by the theorem of Capoyleas
and Pach \cite{CP} (at most $k(2n-2k-1)$ chords with no $k+1$ pairwise crossing, here $k=3$) it has at most $6n-21$
edges. The edges of cyclic length at most $3$ are the $3n-6$ short edges and the six \emph{wrap} edges $(1,n)$,
$(1,n-1)$, $(2,n)$, $(1,n-2)$, $(2,n-1)$, $(3,n)$, which are long. Every long edge of $G$ is among the
$6n-21-(3n-6)=3n-15$ long edges of $H$, so
\[
L\le 3n-15,\qquad e(G)=s+L\le 3n-15+s. \tag{3}
\]
The \emph{deficit} is $d=3n-15-L\ge0$. A graph with $e(G)\ge3n+1$ has $d\le s-16$, so it suffices to rule out
$s\ge16$ with $d\le s-16$.

\subsection{Short edges}
\begin{lemma}\label{lem:len}
Let $F$ be a set of short edges with no non-nested $(r+1)$-matching. Then $F$ has at most $2r$ edges of each length.
If $F$ has $2r$ edges of length $2$, or $2r$ of length $3$, then some $r$ vertices meet every edge of $F$. Hence if no
$r$ vertices meet every edge of $F$, then $|F|\le 6r-2$, with at most $2r-1$ edges of length $2$ and of length $3$.
\end{lemma}

\begin{proof}
The edges of one length form vertex-disjoint paths, and equal-length disjoint edges are never nested; a union of
paths with $m$ edges contains $\lceil m/2\rceil$ disjoint edges. This gives the first bound.

Let $D$ be $2r$ edges of length $\delta\in\{2,3\}$. The bound $\sum\lceil m_P/2\rceil\le r$ over the path components
$P$ forces every component to have an even number of edges. Write each component as $b_0,c_1,b_1,\dots,c_t,b_t$ in
label order and let $C$ be the set of all its vertices $c_j$, over all components; $|C|=r$. An even path keeps a
matching that saturates its $c$-vertices after deleting any one $b$-vertex, or any one edge. Suppose $e\in F$ avoids
$C$. Two $b$-vertices of one component differ by a multiple of $2\delta\ge4$, so the endpoints of $e$ delete at most
one $b$-vertex from each component. If $\delta=2$ no edge of $D$ is nested with $e$, and the remaining $r$-matching of
$D$ extends by $e$. If $\delta=3$ the only possible nested pair is $e=(i+1,i+2)$ inside $f=(i,i+3)\in D$; the
component of $f$ consists of labels $\equiv i\pmod3$ and the endpoints of $e$ have the other two residues, so we may
delete $f$ in its own component and the endpoints of $e$ in theirs, and again find $r$ edges of $D$ disjoint from and
non-nested with $e$. Either way $F$ has a non-nested $(r+1)$-matching, a contradiction. The last statement follows.
\end{proof}

For $r=3$: every $G$ has $s\le18$; if $s\ge17$, three vertices meet every short edge; and if $s=16$ with no such
three vertices, the short edges of lengths $1,2,3$ number $6,5,5$. We say $S$ is \emph{covered by} $C=\{c_1<c_2<c_3\}$
when every short edge meets $C$; then $S$ lies in the union of the \emph{short stars} at the $c_i$ (the at most six
short edges at each $c_i$), and $c_1,c_2,c_3$ are the \emph{centres}.

\begin{lemma}\label{lem:pair}
Let $F$ be a set of short edges, no two of them disjoint and non-nested, and with no common vertex. Then $|F|\le4$,
and $|F|=4$ only for $F=\{ad,bc,x,y\}$ on consecutive vertices $a<b<c<d$ with $x\in\{ab,cd\}$ and $y\in\{ac,bd\}$.
\end{lemma}

\begin{proof}
Four pairwise intersecting edges form a star (an edge meeting $ab$ and $ac$ but not $a$ is $bc$, and no fourth edge
meets all of $ab,ac,bc$). So $F$ has a disjoint pair, which must be nested: $ad,bc$ on consecutive vertices. Every
other edge meets both (otherwise it would be the unique short edge nested with one of them), so lies in
$\{ab,ac,bd,cd\}$, and $ab,cd$ and $ac,bd$ are disjoint non-nested pairs.
\end{proof}

\begin{lemma}\label{lem:uncov}
Let $s=16$ with no three vertices meeting every short edge. Then the six length-$1$ edges are $(c-1,c),(c,c+1)$ for
three \emph{centres} $c_1<c_2<c_3$ with $c_{i+1}-c_i\ge2$, and every short edge that avoids the centres is
$(c-1,c+1)$, $(c-2,c+1)$ or $(c-1,c+2)$ for a centre $c$. If some gap $c_{i+1}-c_i$ is at least $4$, then after
reversing the order of $[n]$ if necessary, $c_2-c_1\ge4$, $c_1\ge4$, and $S$ contains all six short edges at $c_1$.
\end{lemma}

\begin{proof}
By Lemma~\ref{lem:len} there are six length-$1$ edges, with no four disjoint; as there, they form even paths, and the
centres are the second, fourth, $\dots$ vertices of each path. A length-$2$ edge avoiding the centres cannot be nested
with a length-$1$ edge, so deleting its ends must destroy every $3$-matching of length-$1$ edges; this happens only if
both ends are non-centres of one path, that is the edge is $(c-1,c+1)$. For $(u,u+3)$ avoiding the centres, its ends
are not two non-centres of one path (they differ by an odd number), so a $3$-matching of length-$1$ edges avoids them,
and it must contain the one length-$1$ edge nested inside, $(u+1,u+2)$; one of $u+1,u+2$ is a centre.

Assign every short edge to a centre it meets, or to the centre $c$ named above; the edges of the centre $c$ (its
\emph{bucket}) lie in $[c-3,c+3]$. Reversing the order if necessary, suppose $c_2-c_1\ge4$. If $c_3-c_2\ge4$ as
well, each bucket has no two disjoint
non-nested edges: such a pair together with the length-$1$ edges at the other two centres that point away from it is a
non-nested $4$-matching. A bucket with a common vertex has at most six edges, and one without has at most four by
Lemma~\ref{lem:pair}. Three buckets with common vertices would give three vertices meeting every short edge, so the
sizes are $6,6,4$ and two buckets are full short stars at their centres; one of them is at $c_1$ or at $c_3$, and
we reverse the order if it is at $c_3$ (both gaps are at least $4$). If $c_3-c_2\le3$, let $A$ be
the bucket of $c_1$ and $R$ the union of the other two. As before $A$ has no two disjoint non-nested edges and $R$ has
no non-nested $3$-matching (add the length-$1$ edge at $c_1$ pointing left). So $|A|\le6$, and by
Lemma~\ref{lem:len} with $r=2$, $|R|\le12$. If $|R|=12$, then $R$ has four length-$2$ edges and Lemma~\ref{lem:len}
gives two vertices meeting all of $R$, whose full short stars are disjoint, so the two vertices are at least $4$ apart;
but the length-$1$ edges of $R$ are those at $c_2,c_3$, and their only cover by two vertices is $\{c_2,c_3\}$, which
are at most $3$ apart. So $|R|\le11$, $|A|\ge5$, and $A$ has a common vertex, namely $c_1$. If $|A|=5$, then $|R|=11>10$
and Lemma~\ref{lem:len} gives two vertices meeting $R$, which with $c_1$ meet every short edge. Hence $|A|=6$ is the
full short star at $c_1$, and $c_1\ge4$.
\end{proof}

\subsection{Stars of a $3$-triangulation}
Extend $H$ to a maximal set $T$ of chords with no four pairwise crossing, a \emph{$3$-triangulation}. It has exactly
$6n-21$ edges \cite{CP,PS}, among them all $3n-6$ short edges, so it has $3n-15$ long edges. These include all long
edges of $G$, and the set $X$ of long edges of $T$ not in $G$ has exactly $|X|=d$ elements. A \emph{$3$-star} is the
star polygon on seven vertices $v_0,\dots,v_6$ in cyclic order with the seven chords $v_jv_{j+3}$ (indices mod $7$).

\begin{ps}[Pilaud and Santos {\cite[Theorem~1.4(2), Corollary~4.2]{PS}}]
In a $3$-triangulation of the $n$-gon, $n\ge7$, every chord of cyclic length $3$ belongs to exactly one $3$-star, every
chord of cyclic length at least $4$ to exactly two, and the other chords to none.
\end{ps}

Let $1\le i\le n-3$. The chord $(i,i+3)$ has cyclic length $3$, so it lies in a $3$-star of $T$. A chord of a $3$-star
has two star vertices strictly on one side and three on the other, and the short side of $(i,i+3)$ holds only
$i+1,i+2$. So the star has vertices $i,i+1,i+2,i+3,p,q,r$ in cyclic order, with $p,q,r$ in the cyclic interval from
$i+4$ to $i-1$ (first the vertices after $i+3$, then those before $i$), and its other six chords are
\[
(i+1,p),\ (i+2,q),\ (i+3,r),\ (i,p),\ (i+1,q),\ (i+2,r). \tag{4}
\]
Every chord of this star has an endpoint in the \emph{window} $[i,i+3]$ (Figure~\ref{fig:star}). Its chords of length
at least $4$ are long edges of $T$, so they lie in $G$ unless they lie in $X$. Call the window $[i,i+3]$ \emph{good} if,
whatever the positions of $p,q,r$, the set consisting of $S$ and of the chords (4) of length at least $4$ contains a
non-nested $4$-matching. Then:
\begin{quote}
\emph{the star of $T$ on a good window contains an edge of $X$.}
\end{quote}
In particular, if $d=0$ no window is good.

\begin{figure}[ht]
\centering
\begin{tikzpicture}[x=0.8cm,y=0.8cm,v/.style={circle,fill,inner sep=1.2pt}]
\fill[gray!15] (2.7,-0.2) rectangle (6.3,0.2);
\foreach \x/\l in {1/r,3/i,4/i+1,5/i+2,6/i+3,9/p,12/q}
  \node[v,label=below:{\scriptsize$\l$}] at (\x,0) {};
\foreach \x in {0,2,7,8,10,11,13} \node[v,fill=gray] at (\x,0) {};
\draw[gray,thick,dashed] (3,0) to[out=80,in=100] (6,0);
\draw[thick] (4,0) to[out=60,in=120] (9,0);
\draw[thick] (5,0) to[out=60,in=120] (12,0);
\draw[thick] (1,0) to[out=60,in=120] (6,0);
\draw[thick] (3,0) to[out=60,in=120] (9,0);
\draw[thick] (4,0) to[out=60,in=120] (12,0);
\draw[thick] (1,0) to[out=60,in=120] (5,0);
\end{tikzpicture}
\caption{The $3$-star of a $3$-triangulation on the window $[i,i+3]$ (shaded), here with $p,q$ to the right and $r$ to
the left. The dashed chord $(i,i+3)$ is short; the six other chords (4) all have an endpoint in the window. Placed next
to a dense short star, these chords complete a non-nested $4$-matching whatever $p,q,r$ are.}
\label{fig:star}
\end{figure}

\subsection{Good windows}
The next lemma lists the good windows we use. In each part, Table~\ref{tab:windows} gives one non-nested $4$-matching
for each position of $p,q,r$. In every row the four edges are present (an edge of length at most $3$ is one of the
edges assumed to be in $S$, and every other edge listed is a chord (4) whose length the row's inequalities make at least
$4$), they have eight distinct ends, and their left ends and right ends increase together, so no two are nested. The
rows of each part are exhaustive, because a vertex among $p,q,r$ that lies before $i$ is followed in the cyclic order
only by vertices before $i$. Reversing the order of $[n]$ preserves nesting, so each part also holds in mirror image.

\begin{lemma}\label{lem:windows}
Let $c$ be a vertex, and let $A,B$ be disjoint, non-nested edges of $S$. The window $W$ is good in each of the
following cases.
\begin{enumerate}
\item[(a)] $W=[c-2,c+1]$, $(c-3,c),(c,c+2)\in S$, and all ends of $A,B$ are at least $c+3$.
\item[(b)] $W=[c-1,c+2]$, $(c-2,c),(c,c+3)\in S$, and all ends of $A,B$ are at least $c+4$.
\item[(c)] $W=[c-1,c+2]$, $(c-3,c),(c,c+2)\in S$, and all ends of $A,B$ are at least $c+4$.
\item[(d)] $W=[c-1,c+2]$, $(c-3,c),(c,c+3)\in S$, both ends of $A$ are at most $c-4$, and both ends of $B$ are at
least $c+4$ (here $A,B$ need not be non-nested).
\item[(e)] $W=[c-1,c+2]$, $(c-2,c),(c,c+3),(c+1,c+4)\in S$, and both ends of $B$ are at least $c+5$.
\item[(f)] $W=[c-2,c+1]$, $(c-3,c),(c,c+3),(c+2,c+4)\in S$, and both ends of $B$ are at least $c+6$.
\end{enumerate}
\end{lemma}

\begin{table}[ht]
\centering\small
\begin{tabular}{@{}lll@{}}
\toprule
Part & Position of $p<q<r$ (cyclic order) & Non-nested $4$-matching\\
\midrule
(a) & $r\le c-3$ & $(r,c+1),\ (c,c+2),\ A,\ B$\\
    & all right, $r=c+4$ (so $p=c+2$, $q=c+3$) & $(c-3,c),\ (c-2,c+2),\ A,\ B$\\
    & all right, $r\ge c+5$ & $(c-3,c),\ (c-2,p),\ (c-1,q),\ (c+1,r)$\\
\midrule
(b) & $r\le c-2$ & $(r,c+2),\ (c,c+3),\ A,\ B$\\
    & all right, $p=c+3$ & $(c-2,c),\ (c-1,c+3),\ A,\ B$\\
    & all right, $p\ge c+4$ & $(c-2,c),\ (c-1,p),\ (c+1,q),\ (c+2,r)$\\
\midrule
(c) & $p$ right, $p$ before every end of $A,B$ & $(c-1,p),\ (c-3,c),\ A,\ B$\\
    & $r\le c-3$ & $(r,c+1),\ (c,c+2),\ A,\ B$\\
    & $r=c-2$, $q\le c-3$ & $(q,c),\ (c-2,c+2),\ A,\ B$\\
    & $r=c-2$, $p,q$ right, $p$ not before every end & $(c-1,p),\ (c+1,q),\ (c-2,c+2),\ (c-3,c)$\\
    & all right, $p$ not before every end & $(c-1,p),\ (c+1,q),\ (c+2,r),\ (c-3,c)$\\
\midrule
(d) & all right & $A,\ (c-1,p),\ (c,q),\ (c+1,r)$\\
    & $q,r$ left & $(q,c+1),\ (r,c+2),\ (c,c+3),\ B$\\
    & $p,q$ right, $r$ left, $q\ge c+5$ & $A,\ (c-3,c),\ (c-1,p),\ (c+1,q)$\\
    & $p,q$ right, $r$ left, $q=c+4$ (so $p=c+3$) & $A,\ (c-3,c),\ (c-1,c+3),\ B$\\
\midrule
(e) & $r\le c-2$ & $(r,c+2),\ (c,c+3),\ (c+1,c+4),\ B$\\
    & all right, $p=c+3$ & $(c-2,c),\ (c-1,c+3),\ (c+1,c+4),\ B$\\
    & all right, $p\ge c+4$ & $(c-2,c),\ (c-1,p),\ (c+1,q),\ (c+2,r)$\\
\midrule
(f) & $r\le c-3$ & $(r,c+1),\ (c,c+3),\ (c+2,c+4),\ B$\\
    & all right, $r=c+4$ (so $p=c+2$, $q=c+3$) & $(c-3,c),\ (c-1,c+3),\ (c+2,c+4),\ B$\\
    & all right, $r\ge c+5$ & $(c-3,c),\ (c-2,p),\ (c-1,q),\ (c+1,r)$\\
\bottomrule
\end{tabular}
\caption{The certificates of Lemma~\ref{lem:windows}. ``Right'' means after the window and ``left'' before it. In part
(c), if $q=c-3$ the edge $(q,c)$ is the short edge $(c-3,c)$.}
\label{tab:windows}
\end{table}

\begin{proof}
Each row is checked as described above; for instance in the last row of (a), $p\ge c+2$, $q\ge c+3$ and $r\ge c+5$
give lengths $p-c+2$, $q-c+1$, $r-c-1\ge4$, the left ends $c-3<c-2<c-1<c+1$ and the right ends $c<p<q<r$ increase
together. In part (c), the fourth and fifth rows have $q>p\ge c+4$, so $(c+1,q)$ is long. In part (d), the second row
has $q\le c-3$, so $(q,c+1)$ is long. The other rows are similar.
\end{proof}

\subsection{Sixteen short edges}
If $s=16$ then $d=0$ by (3), so it suffices to find a good window. Two finite checks are needed, and they rest on the
following reduction. Fix centres $c_1<c_2<c_3$ for which every short edge lies in $[c-3,c+3]$ for a centre $c$ (true for
covered $S$, and for uncovered $S$ by Lemma~\ref{lem:uncov}). Call a vertex \emph{far} if it is at distance at least
$4$ from every centre; no short edge has a far end or a far vertex strictly inside it.

\begin{lemma}\label{lem:far}
Suppose four consecutive vertices are far, and let $S'$ on $[n-1]$ arise by deleting one of them and relabelling. If
$S'$ has a good window within distance $3$ of a centre, then so does $S$.
\end{lemma}

\begin{proof}
Let $W$ be the corresponding window of $S$; it contains no far vertex. Given $p,q,r$ for $W$, choose one of the four far
vertices outside $\{p,q,r\}$ and delete it. This turns $S$ into $S'$ (no short edge contains it) and keeps the cyclic
order of the star. The non-nested $4$-matching found for $S'$ lifts: its short edges are edges of $S$, a chord of length
at least $4$ after the deletion had length at least $4$ before, and the order of all endpoints is unchanged.
\end{proof}

\begin{lemma}\label{lem:sixteen}
If $s=16$, then some window is good.
\end{lemma}

\begin{proof}
Suppose first that no three vertices meet every short edge, and take the centres of Lemma~\ref{lem:uncov}. If a gap
between consecutive centres is at least $4$, then after reversing the order if necessary $S$ contains all six short
edges at $c_1$ and $c_2\ge c_1+4$, and Lemma~\ref{lem:windows}(c) applies with $c=c_1$, $A=(c_2,c_2+1)$ and
$B=(c_3,c_3+1)$. If both gaps are at most $3$, there are sixteen such families up to translation, on $11$ to $13$
consecutive vertices; this list is part of Computation~1 below. By Lemma~\ref{lem:far} it suffices to place them with
at most seven free vertices on each side, and Computation~1 finds a good window in every case.

Suppose now that $S$ is covered by $C$. Then $S$ is the union of the short stars at the centres with some edges omitted
(and with the stars cut off at $1$ and $n$). If Lemma~\ref{lem:windows}(c), its mirror image, or (d) applies, we are
done; these windows lie within distance $2$ of a centre. Otherwise, by Lemma~\ref{lem:far}, applied repeatedly, it
suffices to treat families with at most three far vertices in each run, that is with gaps $c_{i+1}-c_i\le10$ and at
most three far vertices at each end. Computation~2 lists all of them (with gaps up to $16$ and up to seven free vertices
at each end) and finds a good window within distance $3$ of a centre in each.
\end{proof}

\begin{description}
\item[Computation 1] All sets of $16$ short edges with no non-nested $4$-matching and no three vertices meeting all of
them were enumerated (spans up to $26$); exactly $16$ families up to translation have both centre gaps at most $3$.
Each was placed with $0$ to $7$ free vertices on each side; for every placement and every position of $p,q,r$ the
window $[c-1,c+2]$ at an end centre $c$ carrying all six short edges at $c$ (after reversing the order if needed) is
good. No failures.
\item[Computation 2] All covered families with $s=16$ (gaps $1$ to $16$, each end either cut off or with up to seven
free vertices) were enumerated: $15{,}984$ are not handled by Lemma~\ref{lem:windows}(c) with the extra edge $(c-2,c)$,
by its mirror image, or by (d). Each has a good window $[i,i+3]$ with $c-3\le i\le c$ for a centre $c$. No failures.
\end{description}
Both computations were run by two programs that share no code (see the section on reproducibility).

\subsection{Seventeen and eighteen short edges}
Now $s\in\{17,18\}$, so $S$ is covered by centres $c_1<c_2<c_3$, and it remains to show $|X|\ge s-15$. Let
$s=18-\beta-t-o$, where $\beta$ counts the short edges lost at the two ends ($c_1\le3$ or $c_3\ge n-2$), $t$ the pairs of
centres at distance at most $3$ (whose joining edge is counted twice), and $o$ the omitted short edges at the centres.

At a centre $c$ write $W_c^-=[c-2,c+1]$ and $W_c^+=[c-1,c+2]$, let $T_c^\pm$ be the stars of $T$ on them, and call
$E_c=(c-2,c+2)$ the \emph{ear} at $c$. If $T_c^-=T_c^+$, this star contains the five consecutive vertices
$c-2,\dots,c+2$, hence the chord $E_c$. Two windows that are disjoint cannot lie in one star (that would need eight
vertices).

\begin{lemma}\label{lem:eighteen}
If $s=18$, then $|X|\ge3$.
\end{lemma}

\begin{proof}
Here $\beta=t=o=0$: $S$ consists of the full short stars at the centres, $c_1\ge4$, $c_3\le n-3$ and $c_{i+1}-c_i\ge4$.
All six windows are good: at $c_1$ by Lemma~\ref{lem:windows}(a) and (b) with $A=(c_2,c_2+1)$, $B=(c_3,c_3+1)$; at
$c_3$ by their mirror images with $(c_1-1,c_1)$, $(c_2-1,c_2)$; at $c_2$ by Lemma~\ref{lem:windows}(d) and its mirror
image with $A=(c_1-1,c_1)$ and $B=(c_3,c_3+1)$. Next, the interval strictly between the neighbours of a centre $c$ (or
from $1$, or to $n$) contains no two disjoint non-nested edges of $G$, since such a pair together with the length-$1$
edges at the other two centres pointing away from it would be a non-nested $4$-matching. The edges $E_c$,
$(c-1,c+3)$ and $(c-3,c+1)$ lie in that interval and form such a pair with $(c-3,c)$, $(c-3,c)$ and $(c,c+3)$
respectively, so none of them is in $G$; when one of them is in $T$, it is in $X$.

Suppose $|X|\le2$. Two of the six windows can share a star only if they overlap: $W_c^-$ and $W_c^+$ (an
\emph{internal} identification), or $W_c^+$ and $W_{c'}^-$ for consecutive centres with $c'=c+4$ (a \emph{bridge}). No
three share a star, since any three of the windows span at least eight vertices. A bridge star has exactly the seven
consecutive vertices $c-1,\dots,c+5$, so it contains $(c-1,c+3)$ and $(c+1,c+5)=(c'-3,c'+1)$; both are in $X$, so
$X$ consists of these two, which have no endpoint in a window at the third centre, so by (4) they lie in neither star
there, contradicting goodness. Hence there are no bridges. If $r$ centres have an internal identification, there are
$6-r$ distinct good stars, each containing an edge of $X$, and each edge lies in at most two stars, so $6-r\le2|X|\le4$.
If $r=3$ the three ears are in $X$; so $r=2$ and $X$ consists of the two ears at those centres. The two stars at the
third centre $j$ are distinct and each contains an ear. An ear $E_c$ with $c<c_j$ has its right end $c+2\le c_j-2$, so
by (4) it can lie only in $T^-_{c_j}$, and only if $c_j=c+4$; symmetrically on the right. Hence $j$ is the middle
centre and the centres are $a,a+4,a+8$, with $X=\{(a-2,a+2),(a+6,a+10)\}$.

The star $T^-_{a+4}$ on the window $[a+2,a+5]$ must contain $(a-2,a+2)$, the only edge of $X$ with an end in that
window. In (4) with $i=a+2$ the vertex $i$ has the star neighbours $i+3$ and $p$, so $p=a-2$, and then $q<r$ lie in
$\{a-1,a,a+1\}$. The chords $(a-2,a+3)$, $(q,a+4)$, $(r,a+5)$ have length at least $4$ and are not in $X$, so they are
in $G$; they cross pairwise, and with $(a+8,a+9)$ they form a non-nested $4$-matching. This contradiction proves
$|X|\ge3$.
\end{proof}

\begin{lemma}\label{lem:seventeen}
If $s=17$, then $|X|\ge2$.
\end{lemma}

\begin{proof}
Now $\beta+t+o=1$. We find a centre $c$ whose two windows are good and whose ear is not in $G$, and a good window $J$
disjoint from $[c-2,c+2]$. This suffices: if $T^-_c\ne T^+_c$, the stars $T^-_c$, $T^+_c$ and the star on $J$ are three
distinct stars each containing an edge of $X$, and no edge lies in three stars; if $T^-_c=T^+_c$, this star contains
the ear, which is then in $X$, and if $X$ were just the ear, the star on $J$ would contain it, which is impossible since
the ear has no end in $J$.

If all gaps are at least $4$ ($\beta$ or $o$ is $1$), the defect touches at most one end centre, so the other end
centre $c$ carries all six short edges, and both its windows are good by Lemma~\ref{lem:windows}(a), (b) (or their
mirror images), with $A,B$ the length-$1$ edges at the other two centres pointing away from $c$, or the length-$2$ edge
in the same direction where the defect removes one. Its ear, with $(c-3,c)$ and $A,B$, would be a non-nested
$4$-matching, so it is not in $G$. At the other end centre at least one of the two pairs $\{(c'-3,c'),(c',c'+2)\}$,
$\{(c'-2,c'),(c',c'+3)\}$ (or their mirror images) survives the one defect, so one of its windows is good by the same
lemma; it is disjoint from $[c-2,c+2]$ because the end centres are at least $8$ apart.

Otherwise $t=1$: all three short stars are full and two consecutive centres are close. Reverse the order so that the
close centres are $a$ and $a+h$, $h\in\{1,2,3\}$, and the other centre is $b\ge a+h+4$. Both windows at $b$ are good by
the mirror images of Lemma~\ref{lem:windows}(a), (b), using $(a-2,a),(a-1,a+1)$ if $h=1$ and $(a-1,a),(a+h-1,a+h)$ if
$h=2,3$; and the ear at $b$ is not in $G$ as before. Let $B=(b,b+1)$. For $h=3$ take $J=[a-2,a+1]$, good by
Lemma~\ref{lem:windows}(a) with $c=a$ and $A=(a+3,a+4)$; for $h=2$ take $J=[a-2,a+1]$, good by (f) with $c=a$; for
$h=1$ take $J=[a-1,a+2]$, good by (e) with $c=a$. In each case $J\subseteq[a-2,a+2]$ is disjoint from $[b-2,b+2]$.
\end{proof}

\begin{proof}[Proof of Theorem~\ref{thm:k4plain}]
Suppose $e(G)\ge3n+1$. By (3), $s\ge16$ and $d\le s-16$. If $s=16$ then $d=0$, so $X$ is empty and no window is good,
contradicting Lemma~\ref{lem:sixteen}. If $s=17$ or $18$, then $d=|X|\ge s-15$ by Lemmas~\ref{lem:seventeen}
and~\ref{lem:eighteen}, contradicting $d\le s-16$. Hence $e(G)\le3n$, and with (1), $f_4(n)=3n$.
\end{proof}

\begin{remark}
Lemma~\ref{lem:eighteen} is sharp. For centres with $c_1\ge4$, $c_3\le n-3$ and gaps at least $4$, all $3n-6$ edges
at the three centres together with the six wrap edges form an ordered graph with $3n$ edges, eighteen short edges and
$L=3n-18$. It has no non-nested $4$-matching: a matching contains at most three edges at the centres, and a wrap edge
$(u,v)$ lies in none, since a disjoint edge not nested in it must use one of the at most two vertices outside $[u,v]$.
\end{remark}

\section*{Reproducibility, formal verification and acknowledgement}
All computations, with the scripts and their recorded outputs, are archived with this note (repository folder and
Zenodo record below). For Theorem~\ref{thm:k4plain}: every row of Table~\ref{tab:windows} was checked mechanically
for all positions of $p,q,r$ and all admissible $A,B$ with $n\le 19$, and checks with weakened hypotheses fail, as they
should. Computation~1 was run by two programs, and Computation~2 by a Python program and by an independent program in
C; the two lists of $15{,}984$ families agree. A SAT computation confirms $f_4(n)=3n$ for $8\le n\le13$ (two independent
encodings).

The case analyses of Lemma~\ref{lem:windows}, the exceptional arrangement in the proof of Lemma~\ref{lem:eighteen} and
the counting step $6-r\le2|X|$ are formalised in Lean~4 with Mathlib, as are Theorems~\ref{thm:k3} and~\ref{thm:k4}
(with the theorems of Capoyleas and Pach and of Pilaud and Santos as stated hypotheses); every formalised statement
depends only on the standard axioms, and planted false variants are rejected. The short-edge lemmas, the star
identification and the two computations are not formalised in this version.

The work was done with three AI models, credited by the steps they originated. Codex (OpenAI) found the proof of
Theorem~\ref{thm:k3} (its route, Lemma~\ref{lem:short} and the use of Theorem PS), Lemmas~\ref{lem:len},
\ref{lem:pair} and~\ref{lem:uncov}, parts (a), (b), (e), (f) of Lemma~\ref{lem:windows}, and the proofs of
Lemmas~\ref{lem:eighteen} and~\ref{lem:seventeen}. Claude Opus~5.5 (Anthropic) found the graph $G_9$ and
Theorem~\ref{thm:k4}, the method of single good windows and parts (c), (d) of Lemma~\ref{lem:windows}, the reduction
of the case $s=16$ to Computation~2, ran the computations and the Lean formalisation, and drafted the text. Claude
Fable~5.2 (Anthropic) checked every step independently, supplied the gap in the proof of Lemma~\ref{lem:uncov} for
$|R|=12$, the weakening of part (c), and Lemma~\ref{lem:far}, and wrote the independent programs. The author is
responsible for the content.

\begin{thebibliography}{9}
\bibitem{BFT} J.~Bar\'at, A.~Freschi and G.~T\'oth, Matchings avoiding ordered patterns, arXiv:2512.15461v2 (2025).
\bibitem{BGT} J.~Bar\'at, A.~Gy\'arf\'as and G.~T\'oth, Monochromatic spanning trees and matchings in ordered
complete graphs, \emph{J. Graph Theory} 105(4) (2024), 523--541, doi:10.1002/jgt.23058.
\bibitem{CC33} Open problems of the 33rd Workshop on Cycles and Colourings, arXiv:2511.02892 (2025), Problem 1.1
(J.~Bar\'at).
\bibitem{CP} V.~Capoyleas and J.~Pach, A Tur\'an-type theorem on chords of a convex polygon, \emph{J. Combin. Theory
Ser. B} 56 (1992), 9--15, doi:10.1016/0095-8956(92)90003-G.
\bibitem{dC} P.~M.~M.~de Castro, Minimum blockers and extremal graphs for nonnested matchings, arXiv:2609.14484v2
(2026).
\bibitem{PS} V.~Pilaud and F.~Santos, Multitriangulations as complexes of star polygons, \emph{Discrete Comput.
Geom.} 41 (2009), 284--317, doi:10.1007/s00454-008-9078-6; arXiv:0706.3121 (numbering as in the arXiv version).
\end{thebibliography}
\end{document}
